% Suggested packages for the main Overleaf document:
% \usepackage{graphicx}
% \usepackage{tikz}
% \usetikzlibrary{positioning,arrows.meta,fit,backgrounds}
% \usepackage{xcolor}

\section{Task Decomposition of the Prism Adaptation Workflow}
\label{sec:task-decomposition}

This section visualises how the experiment is structured from the participant's point of view. Rather than treating the implementation as a collection of scripts and interfaces, the design is described as a small set of repeated interaction loops. The overall workflow follows a \emph{measurement--exposure--measurement} structure: a baseline task is first performed without perturbation, then an exposure block is completed under one of the implemented visual effects, and finally the same measurement task is repeated in a post-exposure phase. The key distinction is therefore between tasks that \emph{measure} a spatial estimate and the task that is intended to \emph{change} it through repeated visuomotor correction.

An additional design consideration concerns how responses are confirmed and what visual information is available before and after an action. The default controller implementation uses trigger confirmation, whereas embodied interaction uses dwell-based confirmation. The system also includes prototype-specific Heisenberg-mitigation variants, such as opposite-hand confirmation and controller-hover confirmation. These variants should therefore be understood as configurable design alternatives in the prototype rather than as fixed assumptions inherited from the prism-adaptation literature. Across the tasks, the most important feedforward cues are the currently shown target or stimulus layout, while the most important feedback cues are the accepted response, the loading/progress indicator, and the visible shift to the next required interaction. Figure~\ref{fig:prism-task-cycles} therefore shows the abstract task structure, while a separate movement illustration can be used alongside it to show the concrete body posture and arm motion associated with each task.

\begin{figure}[t]
    \centering
    \resizebox{\linewidth}{!}{%
    \begin{tikzpicture}[
        x=1cm,
        y=1cm,
        >=Latex,
        titlebox/.style={draw, rounded corners=2pt, align=center, fill=blue!8, minimum width=3.8cm, minimum height=0.85cm, font=\small\bfseries},
        taskbox/.style={draw, rounded corners=2pt, align=center, fill=blue!6, minimum width=2.8cm, minimum height=0.8cm, font=\small\bfseries},
        subbox/.style={draw, rounded corners=2pt, align=center, fill=green!8, minimum width=2.9cm, minimum height=0.8cm, font=\small},
        actionline/.style={draw, rounded corners=2pt, align=center, fill=yellow!12, minimum width=10.2cm, minimum height=0.75cm, font=\scriptsize, text width=9.8cm},
        noteff/.style={draw, rounded corners=2pt, align=left, fill=green!5, minimum width=6.0cm, minimum height=1.15cm, font=\scriptsize, text width=5.6cm},
        notefb/.style={draw, rounded corners=2pt, align=left, fill=blue!5, minimum width=6.0cm, minimum height=1.15cm, font=\scriptsize, text width=5.6cm},
        arr/.style={-Latex, line width=0.7pt}
    ]
        % Workflow
        \node[titlebox] at (2.2,0) (base) {Baseline measurement};
        \node[titlebox] at (6.7,0) (exp) {Exposure under effect};
        \node[titlebox] at (11.2,0) (post) {Post measurement};
        \draw[arr] (4.15,0) -- (4.75,0);
        \draw[arr] (8.65,0) -- (9.25,0);

        % Open-loop / line bisection
        \node[taskbox, anchor=west] at (0,-2.2) (t1) {Open-loop\\Line bisection};
        \node[subbox] at (4.5,-2.2) (s11) {Reset posture};
        \node[subbox] at (8.2,-2.2) (s12) {Estimate and aim};
        \node[subbox] at (11.9,-2.2) (s13) {Confirm and reset};
        \draw[arr] (6.00,-2.2) -- (6.65,-2.2);
        \draw[arr] (9.70,-2.2) -- (10.35,-2.2);

        \node[actionline] (actmeas) at (8.2,-3.55) {Arms lowered $\rightarrow$ raise arm into frontal workspace $\rightarrow$ aim at target or midpoint $\rightarrow$ trigger/dwell $\rightarrow$ lower arm};

        \node[noteff] (ffmeas) at (4.7,-5.25) {\textbf{Feedforward:} visible target or line, board plane, and the appearance of the next measurement layout.};
        \node[notefb] (fbmeas) at (11.5,-5.25) {\textbf{Feedback:} signed readout after acceptance and a short loading/progress bar before the next trial.};
        \node[draw, dashed, rounded corners=3pt, fit=(t1)(s11)(s12)(s13)(actmeas)(ffmeas)(fbmeas), inner sep=6pt] (g1) {};
        \node[font=\scriptsize\bfseries, anchor=west, fill=white, inner sep=1pt] at ([xshift=2pt,yshift=2pt]g1.north west) {Measurement block};

        % Landmark
        \node[taskbox, anchor=west] at (0,-7.4) (t2) {Landmark};
        \node[subbox] at (4.5,-7.4) (s21) {Inspect stimulus};
        \node[subbox] at (8.2,-7.4) (s22) {Compare sides};
        \node[subbox] at (11.9,-7.4) (s23) {Select response};
        \draw[arr] (6.00,-7.4) -- (6.65,-7.4);
        \draw[arr] (9.70,-7.4) -- (10.35,-7.4);

        \node[actionline] (actland) at (8.2,-8.75) {Inspect line $\rightarrow$ compare left and right segments $\rightarrow$ choose side $\rightarrow$ confirm};

        \node[noteff] (ffland) at (4.7,-10.45) {\textbf{Feedforward:} the visible landmark stimulus and the response side it affords.};
        \node[notefb] (fbland) at (11.5,-10.45) {\textbf{Feedback:} accepted choice and the loading bar before the next trial.};
        \node[draw, dashed, rounded corners=3pt, fit=(t2)(s21)(s22)(s23)(actland)(ffland)(fbland), inner sep=6pt] (g2) {};
        \node[font=\scriptsize\bfseries, anchor=west, fill=white, inner sep=1pt] at ([xshift=2pt,yshift=2pt]g2.north west) {Assessment block};

        % Exposure
        \node[taskbox, anchor=west] at (0,-12.6) (t3) {Exposure};
        \node[subbox] at (4.5,-12.6) (s31) {Detect target shift};
        \node[subbox] at (8.2,-12.6) (s32) {Acquire and stabilise};
        \node[subbox] at (11.9,-12.6) (s33) {Confirm and retarget};
        \draw[arr] (6.00,-12.6) -- (6.65,-12.6);
        \draw[arr] (9.70,-12.6) -- (10.35,-12.6);

        \node[actionline] (actexp) at (8.2,-13.95) {See active target $\rightarrow$ redirect aim $\rightarrow$ hold or trigger $\rightarrow$ hit/miss outcome $\rightarrow$ continue to next target};

        \node[noteff] (ffe) at (4.7,-15.65) {\textbf{Feedforward:} active-target highlighting, fixed left--center--right layout, and the loading/layout shift when the next target appears.};
        \node[notefb] (fbe) at (11.5,-15.65) {\textbf{Feedback:} hit or miss outcome, hit marker, and visible shift to the next required target.};
        \node[draw, dashed, rounded corners=3pt, fit=(t3)(s31)(s32)(s33)(actexp)(ffe)(fbe), inner sep=6pt] (g3) {};
        \node[font=\scriptsize\bfseries, anchor=west, fill=white, inner sep=1pt] at ([xshift=2pt,yshift=2pt]g3.north west) {Exposure block};
    \end{tikzpicture}}
    \caption{Visual task decomposition of the implemented workflow. Open-loop pointing and line bisection share a reset--aim--confirm structure, landmark is a comparative judgment task, and exposure is a repeated target-to-target acquisition task under perturbation.}
    \label{fig:prism-task-cycles}
\end{figure}

\subsection{Open-Loop Pointing}
\label{subsec:open-loop-decomposition}

Open-loop pointing is the clearest measurement task in the system because each trial yields a single signed endpoint error relative to one visible target. The participant begins with the arm lowered alongside the body, raises it into the frontal workspace in a semi-straight configuration, aims toward the target, commits the response, and lowers the arm again. In controller mode, commitment occurs through the trigger; in embodied mode, the endpoint is accepted once the fingertip remains stable for approximately 1.5\,s. The most important feedforward cue is simply the visible target location on the board, while the most important feedback is the accepted endpoint together with the short loading bar and transition to the next trial.

\subsection{Line Bisection}
\label{subsec:line-bisection-decomposition}

Line bisection uses almost the same motor pattern as open-loop pointing, but the participant is no longer given a point target. Instead, a horizontal line is shown and the participant must infer its midpoint before aiming. The task can therefore be understood as the same reset--aim--confirm loop combined with a stronger perceptual judgment component. The visible line provides feedforward about where the midpoint should be estimated, while the main feedback again consists of the accepted response, the resulting readout, and the transition to the next trial. The output measure is the signed deviation from the true midpoint.

\subsection{Landmark}
\label{subsec:landmark-decomposition}

Landmark has a different structure from the other tasks because it is not primarily a target-acquisition task. The participant inspects the displayed stimulus, compares the two visible line segments, decides which side appears longer, and then selects the corresponding response. For that reason, landmark is better described as a comparative perceptual judgment followed by a discrete choice. Its feedforward is the visual organisation of the stimulus itself, while its feedback is the confirmation that the choice has been accepted and that the next trial is loading. In the present project, landmark is retained mainly as a supplementary assessment task.

\subsection{Exposure}
\label{subsec:exposure-decomposition}

Exposure is the adaptation-inducing task and is therefore structured differently from the measurement tasks. Instead of returning to a start posture between trials, the participant continuously redirects the response between three horizontally arranged targets: left, center, and right. One target is active at a time, and after each confirmed attempt the next target appears. In controller mode, confirmation is normally trigger-based, although opposite-hand confirmation and controller-hover confirmation are also supported; in embodied mode, confirmation is dwell-based. The dominant feedforward cues are the active-target highlight, the fixed three-target layout, and the loading or transition cue that precedes the next required movement. The dominant feedback cues are the hit or miss outcome, the hit marker, and the visible shift to the next target under the active perturbation.

\subsection{Measurement Tasks Versus Exposure}
\label{subsec:measurement-vs-exposure}

The workflow therefore contains two different experimental logics. Open-loop pointing and line bisection are measurement tasks that isolate a single accepted response and compare it between baseline and post-exposure phases. Exposure, by contrast, is a repeated target-to-target acquisition task whose purpose is to induce sensorimotor adaptation under perturbation. Landmark remains conceptually separate as a perceptual comparison task that may be useful for later clinical use, but is less central to the present healthy-participant workflow.
